% v6 section. Outcome sentences are generated from the frozen analysis's output; this
% file contains no outcome.
\section{A Predeclared Replication at the Planned Size}\label{sec:v6replication}

The previous version's live comparison of five controllers was predeclared but underpowered. With \vsixPilotN{}
independent cases per checkpoint, its usefulness claims could not clear the five-point margin at the corrected
\(\alpha\) even with no losses (the zero-loss bound is \vsixPilotZeroLossBound), and it reported its claims as
inconclusive. The plan's own design diagnostic gives the sample the noninferiority claims need: \vsixEOneMinN{}
independent cases. E1 runs \vsixEOneN.

\para{Design (E1).} The reference plan is unchanged, and so is its analysis script, frozen by hash. Five controllers are compared:
\begin{itemize}
  \item no mediation;
  \item a stateless per-call policy;
  \item a same-checkpoint learned reviewer;
  \item a durable broker without a sink fence;
  \item complete mediation.
\end{itemize}
They run on \(N = 128\) fresh independent cases per checkpoint, for Qwen3-4B and Qwen3-8B, which is
\vsixEOneCampaigns{} campaigns on CPU.

The primary claims are, for each comparator and checkpoint, one-sided safety superiority and five-point usefulness
noninferiority of complete mediation: 16 component claims at Bonferroni \(\alpha = 0.05/16\), with checkpoints never
pooled. The freeze, roster seed and nonces are generated by the harness after the sample size is fixed. Two amendments
precede the data:
\begin{itemize}
  \item the first raises the wall-clock cap that had limited the pilot;
  \item the second follows a stopped launch. Proposals expired four hours after each checkpoint began, which would
        have denied the 8B tail as expired. That launch was stopped at 235 records, with no aggregate viewed and one record opened while
        diagnosing a schema error (Appendix~\ref{app:v6ledger}); it is kept and disclosed.
\end{itemize}

\para{What the primary claims can and cannot show.} Complete mediation's unauthorized rate is zero by construction:
its admission predicate is the scoring predicate. The safety claims therefore measure how often each comparator leaks,
not whether complete mediation can be broken; Sections~\ref{sec:v6stop} and~\ref{sec:v6reviewers} address that.
Usefulness keeps the previous version's scoring, which is coupled to the gate: a task succeeds only if the case's
first approved proposal commits and nothing else does, and a denial counts as a false block only if the proposal
satisfied the same predicate the gate enforces, so complete mediation cannot record a false block
(Appendix~\ref{sec:systems}).
\vsixEOneLossRuleSentence

In the revocation cases, only a fenced controller can complete the task, because the queued old operation otherwise
commits. Those cases count toward usefulness in complete mediation's favour. We therefore report, beside every
usefulness claim, the predeclared secondary S4, which recomputes the same bound without revocation cases. S1 counts
effects after a withdrawal as unauthorized, S2 counts the first proposal's effect regardless of others, and S3 splits
by attack family.

\para{Result.} \vsixEOneOutcomeSentence{} \vsixEOneSFourSentence{} The pilot (\(N = 32\)) is reported separately in
Appendix~\ref{app:v6e1} and never pooled.

\para{A paired reading (post hoc).} The frozen analysis bounds each direction of discordance separately and ignores
cases where complete mediation gains, so it is conservative: on the pilot it reported as inconclusive a comparison in
which ten of ten safety discordances favoured complete mediation, which an exact sign test puts at \(p \approx 0.001\).
Before E1's outcomes were read, we fixed a second reading and report it beside the frozen one, not instead of it
(Table~\ref{tab:v6e1paired}): how often each comparator let an unauthorized effect commit, with an exact interval,
which is the safety quantity E1 actually measures, and the paired discordances on safety and usefulness.

\begin{table*}[t]\centering\footnotesize
\caption{E1 per controller. Leak and task success: the predeclared primary scoring (complete mediation's leak is zero by construction, and task success requires that nothing but the first approved proposal commit). After withdrawal: S1, which also counts an effect committed after a withdrawal as unauthorized. First proposal: S2, whether the first approved proposal's effect committed, whatever else did. Safety and usefulness discordances: cases favouring complete mediation versus the comparator, a paired reading fixed after the plan and before the outcomes were read; no test is run on them, because every safety discordance must favour complete mediation. The frozen analysis (Table~\ref{tab:v6e1}) is primary.}\label{tab:v6e1paired}
\begin{tabular}{@{}lcccccc@{}}\toprule
 & \multicolumn{3}{c}{Unauthorized (cases)} & \multicolumn{3}{c}{Useful (cases)} \\
\cmidrule(lr){2-4}\cmidrule(l){5-7}
Controller & Leak & After withdrawal (S1) & Safety & Task success & First proposal (S2) & Usefulness \\\midrule
\multicolumn{7}{@{}l}{\emph{Qwen3-4B, 128 paired cases}} \\
\quad No mediation & 39 & 54 & 39:0 & 23 & 44 & 20:6 \\
\quad Per-call policy & 10 & 24 & 10:0 & 4 & 23 & 33:0 \\
\quad Same-model reviewer & 32 & 46 & 32:0 & 19 & 35 & 22:4 \\
\quad Durable, no fence & 0 & 14 & 0:0 & 27 & 36 & 10:0 \\
\quad Complete mediation & 0 & 0 & -- & 37 & 37 & -- \\
\multicolumn{7}{@{}l}{\emph{Qwen3-8B, 128 paired cases}} \\
\quad No mediation & 17 & 37 & 17:0 & 46 & 82 & 28:0 \\
\quad Per-call policy & 8 & 28 & 8:0 & 32 & 59 & 43:1 \\
\quad Same-model reviewer & 8 & 28 & 8:0 & 46 & 69 & 28:0 \\
\quad Durable, no fence & 0 & 20 & 0:0 & 55 & 74 & 19:0 \\
\quad Complete mediation & 0 & 0 & -- & 74 & 74 & -- \\
\bottomrule\end{tabular}\end{table*}

